% Where each model combines time. Orange marks the fusion point.
\begin{figure}[t]
\centering
\begin{tikzpicture}[
    font=\scriptsize, >={Stealth[length=3pt]}, node distance=3.5mm,
    box/.style={draw, rounded corners=1pt, fill=black!6, minimum height=4.2mm,
                inner sep=2pt, align=center},
    fuse/.style={box, fill=orange!35},
    lab/.style={anchor=west, font=\scriptsize\bfseries, minimum width=12mm, inner sep=0},
    io/.style={inner sep=1pt, align=center},
]
% Per-frame aggregation
\node[lab] (l1) at (0, 0) {Per-frame};
\node[io, right=1mm of l1] (x1) {$x_t$};
\node[box, right=of x1] (c1) {2D CNN + FC};
\node[box, right=of c1] (s1) {softmax};
\node[fuse, right=of s1] (f1) {mean$_t$};
\node[io, right=of f1] (y1) {$\hat y$};
\draw[->] (x1) -- (c1); \draw[->] (c1) -- (s1); \draw[->] (s1) -- (f1); \draw[->] (f1) -- (y1);
% Late fusion
\node[lab] (l2) at (0, -0.7) {Late};
\node[io, right=1mm of l2] (x2) {$x_t$};
\node[box, right=of x2] (c2) {2D CNN};
\node[fuse, right=of c2] (f2) {mean$_t\,h_t$};
\node[box, right=of f2] (s2) {FC};
\node[io, right=of s2] (y2) {$\hat y$};
\draw[->] (x2) -- (c2); \draw[->] (c2) -- (f2); \draw[->] (f2) -- (s2); \draw[->] (s2) -- (y2);
% Early fusion
\node[lab] (l3) at (0, -1.4) {Early};
\node[io, right=1mm of l3] (x3) {$x_{1:T}$};
\node[fuse, right=of x3] (f3) {stack: $3T$ ch.};
\node[box, right=of f3] (c3) {2D CNN + FC};
\node[io, right=of c3] (y3) {$\hat y$};
\draw[->] (x3) -- (f3); \draw[->] (f3) -- (c3); \draw[->] (c3) -- (y3);
% 3D CNN
\node[lab] (l4) at (0, -2.1) {3D CNN};
\node[io, right=1mm of l4] (x4) {$x_{1:T}$};
\node[fuse, right=of x4] (c4) {3D CNN ($3{\times}3{\times}3$ kernels)};
\node[box, right=of c4] (s4) {FC};
\node[io, right=of s4] (y4) {$\hat y$};
\draw[->] (x4) -- (c4); \draw[->] (c4) -- (s4); \draw[->] (s4) -- (y4);
% Two-stream
\node[lab] (l5) at (0, -3.0) {Two-stream};
\node[io, right=1mm of l5, yshift=2.6mm] (x5) {$x_{1:T}$};
\node[box, right=of x5] (c5) {late fusion};
\node[io, right=1mm of l5, yshift=-2.6mm] (u5) {$u_{1:T-1}$};
\node[fuse, right=of u5] (g5) {stack: $3(T{-}1)$ ch.\ + CNN};
\node[fuse, right=of g5, yshift=2.6mm] (f5) {$(1{-}w)\,p_\mathrm{rgb}$\\$+\,w\,p_\mathrm{flow}$};
\draw[->] (x5) -- (c5); \draw[->] (u5) -- (g5);
\draw[->] (c5.east) -- node[above, pos=0.5] {$p_\mathrm{rgb}$} (c5.east -| f5.west);
\draw[->] (g5.east) -- ++(1.5mm, 0) |- ([yshift=-1mm]f5.west);
\end{tikzpicture}
\caption{Where each model combines the $T{=}10$ frames (orange). All models
share the same 2D backbone ($x_t$: RGB frame, $h_t$: its feature vector,
$u_t$: colour-coded RAFT flow from frame $t$ to $t{+}1$). The two streams are
trained separately and fused by weighting their softmax scores.}
\label{fig:models}
\end{figure}


\section{Models}
\label{sec:method}
\paragraph{Shared backbone.}
Every variant uses the same VGG-style~\cite{simonyan2015vgg} network: four
stages, each of two $3{\times}3$
convolution--batch-norm~\cite{ioffe2015batchnorm}--ReLU layers
(64/128/256/256 channels) followed by $2{\times}2$ max pooling. Global
average pooling gives a frame feature $h(x_t)\in\mathbb{R}^{256}$, and
dropout (0.5) precedes a linear classifier (2.33\,M parameters).

\paragraph{Why this backbone.}
To measure fusion rather than optimisation, every variant must be able to
fit its training set. Our first backbone, three convolutions without
normalisation, failed this check. Per-frame aggregation reached 90.8\%
validation accuracy, but early fusion and the 3D CNN fitted only 29\% and
42\% of their training videos. Their low validation accuracy (22.5\% and
32.5\%) therefore reflected underfitting, not missing temporal information.
The current backbone adds batch normalisation, which stabilises
optimisation, and a fourth stage, which enlarges the receptive field from
22\,px (10\% of the frame) to 76\,px (68\%). Together with the warm starts
below, these changes let every variant fit its training set
(\cref{sec:experiments}).

\paragraph{Fusion variants (\cref{fig:models}).}
With logits $z_t=Wh(x_t)+b$ for frame $t$ and $T=10$:
\begin{align}
  p_\mathrm{frame}(y\mid V) &= \tfrac{1}{T}\textstyle\sum_t \operatorname{softmax}(z_t)_y,\\
  p_\mathrm{late}(y\mid V)  &= \operatorname{softmax}\big(W\,\tfrac{1}{T}\textstyle\sum_t h(x_t)+b\big)_y.
\end{align}
The per-frame model trains on single frames with their video's label;
late fusion trains on whole videos. Both ignore frame order. Early fusion
stacks the frames into a 30-channel image, so only its first convolution
is wider (2.35\,M parameters) and it sees frame order from the first layer.
The 3D CNN uses the same layer layout with $3{\times}3{\times}3$ kernels;
pooling keeps time in the first two stages and halves it in the last two
($T{:}\;10\to10\to5\to2$), so its final features span the whole clip
(6.98\,M parameters).

\paragraph{Warm starts.}
Training video models from scratch on 500 videos is difficult, so late
fusion, early fusion, and the 3D CNN start from the trained per-frame
model of the same split. Late fusion copies its weights. Early fusion
repeats the RGB input kernel over the $T$ frames and divides it by
$T$~\cite{wang2016tsn}. The 3D CNN inflates every kernel,
$W^{3D}_{:,:,\tau}=W^{2D}/3$~\cite{carreira2017i3d}. For a static clip,
both inflated models initially reproduce the per-frame features (approximately, for the 3D
CNN), so any
gain must come from motion learned during training.

\paragraph{Two streams.}
The temporal stream is an early-fusion network over the nine stacked flow
fields (27 input channels)~\cite{simonyan2014twostream}. The RGB stream is
the late-fusion model trained on the corrected splits. As in the original
two-stream work, the streams are trained separately and combined as
$p=(1-w)\,p_\mathrm{rgb}+w\,p_\mathrm{flow}$, with $w\in\{0,0.05,\dots,1\}$
chosen on validation; we also report $w=0.5$. As an alternative, a jointly
trained dual-stream model concatenates both streams' features before one
linear layer (4.67\,M parameters), with its RGB encoder initialised from
the per-frame model.
